\paragraph{Implementation}

The numerical implementation follows the tensor-reconstruction strategy
introduced above for the quartic case and extends it to sextic
centrifugal distortion.
Spectroscopic parameters are first mapped onto the corresponding tensor
coefficients through Moore--Penrose pseudoinverse reconstruction.
Representation changes are then performed directly in tensor space by
permutation of the tensor indices corresponding to the chosen
relabeling of the principal axes.

The resulting computational workflow consists of three steps:

\begin{enumerate}
\item reconstruction of the tensor coefficients from the spectroscopic
constants through pseudoinverse reconstruction,
\item permutation of the tensor indices according to the desired
principal-axis transformation,
\item recomputation of the spectroscopic constants in the new
representation by explicit forward projection.
\end{enumerate}

This procedure applies uniformly to both quartic and sextic centrifugal
distortion constants. Once the tensor-reconstruction machinery has been
established, representation transforms no longer require
representation-specific analytical transformation matrices.

For the quartic problem the implementation works directly with the
reduced symmetric tensor $\tau'$, together with its spectral
invariants and, when needed, with the practical $3+2$ coordinates of
the pseudoinverse slice. For the sextic problem the implementation
uses the canonical decomposition
\[
\mathbf H_S=\mathbf B_{\mathrm{rep}}\boldsymbol\sigma+\mathbf r,
\]
in which $\boldsymbol\sigma$ spans the physical five-dimensional
sextic sector and $\mathbf r$ measures the complementary residual.
The same five-dimensional content may be interpreted structurally as
one isotropic scalar plus four amplitudes in the even harmonic
degree-six sector, or, in the representation-adapted viewpoint
relevant to CeDiTT, as a $1\oplus1\oplus1\oplus2$ decomposition under
cyclic relabelings of the principal axes.

In addition to performing representation transforms, the program also
evaluates the parameter $s_{111}$ introduced above. Within the present
framework $s_{111}$ is retained as a diagnostic quantity associated
with the chosen Watson parametrization rather than as a fundamental
tensor coordinate. It therefore remains useful for practical analysis
and numerical monitoring even though the tensor formulation itself is
expressed directly in terms of reduced quartic and sextic tensors.


\paragraph{Invariant--gauge analysis}

The tensor formulation provides a complementary perspective on
centrifugal distortion parameters. Watson constants are convenient
coordinates for spectral analysis, but they do not constitute the
primary tensorial objects. In the quartic case the natural reduced
tensor is $\tau'$, while in the sextic case the corresponding physical
content is represented by the canonical five-dimensional sextic sector
described above.

From this viewpoint the inverse reconstruction problem naturally takes
the form of an affine gauge problem: spectroscopic constants determine
equivalence classes of tensor representatives, and the Moore--Penrose
pseudoinverse selects a specific representative within each class.
The associated gauge fixing is therefore a property of the
reconstruction algorithm rather than a physical constraint on
centrifugal distortion.

For this reason invariant analysis is most naturally performed at the
tensor level rather than directly in terms of Watson's constants.
In the quartic case this means working with the spectral invariants of
$\tau'$ and with the corresponding tensor representatives on the
pseudoinverse slice. In the sextic case it means working with the
canonical physical coordinates $\boldsymbol\sigma$, the residual
vector $\mathbf r$, and the equivalent harmonic or
representation-adapted decompositions of the same five-dimensional
sector.

The tensor representation therefore provides a common framework within
which rovibrational information obtained from spectroscopic fits,
direct quantum--chemical calculations, or rovibrational perturbation
theory can be compared on the same footing. At the same time, the
reconstructed tensor can always be projected back onto a chosen
Watson reduction and principal-axis representation whenever a
spectroscopic parametrization is required.

The overall structure of the tensor-based framework implemented in the
present program is illustrated in Fig.~\ref{fig:tensor_framework}.
Different sources of rovibrational information are first mapped onto a
common tensor representation, which can subsequently be used either for
invariant analysis or for transformations between principal-axis
representations.

In this sense the tensor reconstruction step provides the conceptual
and computational bridge between the direct quantum-mechanical
evaluation of rovibrational couplings and the inverse spectroscopic
determination of Watson parameters.
